\section{마무리 (30분) — PM 트랙을 마치며 · Product 트랙 예고}

% =====================================================================
% 마무리 — 교재 제6장 · 워크북 §7
% =====================================================================
\begin{frame}{이 교시에서 — 30분}
\begin{itemize}
\item 학술 배경 6갈래 — 표준이 "하라"고만 말하는 곳의 "왜"와 "언제 틀리는가"
\item 읽을 자료 — 필수 4 / 심화 8 / 참고 / 읽지 않아도 되는 것
\item 4일 20종 산출물의 사슬을 거꾸로 — ⑳에서 ①로 (그림 6-1)
\item 세 표준의 전환기 — PMBOK 8 · ITIL v5 · 개인정보 보호법 제21445호
\item Day1\til{}4 네 문장 · Product 트랙 예고 · 사후 읽기와 30일 과제
\end{itemize}
\keymsg{오늘의 한 문장 — 오픈은 이벤트가 아니라 이관이며, 이관은 받는 쪽이 서명해야 끝난다}
\note{교재 제6장과 "PM 트랙을 마치며", 워크북 §7. 6교시 실습 ②(이관·종료·교훈·원장)를 막 끝낸 상태이므로, 이 교시는 새 내용을 쌓지 않고 4일을 하나의 사슬로 묶는다. 학술 배경은 문헌 소개가 아니라 각 문헌의 "흔한 오독" 한 줄에 집중한다 — 오독을 말해야 실무자가 기억한다. 후반 15분은 Product 트랙 예고와 30일 과제에 쓴다. 예상 소요 2분.}
\end{frame}

\begin{frame}{학술 배경 I — 책임: 누가 편익에 서명하는가}
\begin{tbl}[\scriptsize]{L{2.6cm}L{4.2cm}L{4.6cm}}
문헌 & 이 교재가 빌린 것 & 흔한 오독 \\ \midrule
Zwikael \& Smyrk (2011) & ITO 모형 · output/outcome 책임 분리 · 프로젝트 소유자 \ra{} ⑳ ITO 서명표 & "PM은 편익에 무관심해도 된다" — 아니다, utilisation이 가능한 output을 설계할 책임(두 분모 리포트 Must) \\
Ward \& Daniel (2012) & 편익 의존성 네트워크 · 편익 소유자 = 업무 변화 소유자 \ra{} 5.5절 세 실패 원인 & 편익 관리 = 종료 후 측정 — 아니다, 착수 때 네트워크를 그려 업무 변화 없는 편익을 미리 드러내는 것 \\
Williams (2008) · Duffield \& Whitty (2015) & 교훈은 수집되나 적용되지 않는다는 실증 · Syllk 6층 \ra{} ⑲ "적용 대상·소유자" 칸 & Syllk = 지식관리 도구 도입 논거 — 아니다, technology는 6층 중 하나, process·infrastructure 없으면 도구는 저장소 \\
\end{tbl}
\keymsg{outcome 소유자가 없는 프로젝트는 시작하면 안 된다 — 저자들의 실제 주장이다}
\note{교재 6.1\til{}6.3. 세 문헌은 모두 "책임"에 관한 것이다. Zwikael \& Smyrk에서 빌린 규칙은 "output 서명은 G4에서 종결된다"이고, 그래서 원장의 ⑤ 칸이 서명란이다(그림 5-2). Ward \& Daniel의 초점이 착수 시점이라는 점은 Day1 BC §4의 "측정 책임자" 열과 연결해 말한다 — 그 열이 부서명이었다면 오늘 서명할 사람이 없었다. Williams의 게재지는 리딩리스트 표기와 달라 [검증 필요]로 남겨 두었음을 정직하게 말한다. 예상 소요 4분.}
\end{frame}

\begin{frame}{학술 배경 II — 기억과 속도: 왜 교훈은 남지 않고, 안정은 환전되는가}
\begin{tbl}[\scriptsize]{L{2.6cm}L{4.2cm}L{4.6cm}}
문헌 & 이 교재가 빌린 것 & 흔한 오독 \\ \midrule
Lundin \& Söderholm (1995) & 임시조직의 Transition · 제도화된 종료 \ra{} 종료 직전 이탈과 교훈 실패의 구조적 설명 & 종료를 "행사"로 읽는 것 — 처방은 종료의 제도화, 즉 오늘의 다섯 문서 \\
Beyer 외, SRE (2016) & error budget = "안정과 속도의 환율" · SLI/SLO/SLA 구분 \ra{} 219.2분 · 릴리스 몫 87.7 & "다운을 허용하는 면죄부" — 아니다, 릴리스 동결 규칙(policy)을 숫자로 갖는 것 \\
Ford, Ford \& D'Amelio (2008) & 저항은 변화 설계에 관한 정보, 때로 추진자가 만든 것 \ra{} 3.9절 ADKAR 장벽 격자 & 저항 = 극복 대상 — 아니다, 가맹 K/A · 물류 A/D 장벽은 설계의 피드백 \\
\end{tbl}
\keymsg{함께 읽을 것 — Accelerate(변경 실패율·복구 시간)는 환전의 입력, ADKAR·Kotter·Bridges는 저항의 진단 도구}
\note{교재 6.4\til{}6.6. Lundin \& Söderholm은 Day1 1.3절에서 네 구조 개념을 빌렸던 논문이고, 오늘 Transition을 회수한다 — 4일이 한 논문의 앞뒤로 닫히는 셈이다. SRE는 3.6절의 계산 원전이며, 3장·4장은 온라인 공개라 필수 자료에 넣었다. Ford 외는 3.9절의 태도를 정한 논문이다 — 마감 22:00 저항(R)을 "가맹점주가 틀렸다"로 적지 않고 이관 3번(승인 전)으로 적은 이유가 여기 있다. 예상 소요 4분.}
\end{frame}

\begin{frame}{읽을 자료 — 필수(과정 전) · 심화(과정 후) · 읽지 않아도 되는 것}
\begin{tbl}[\scriptsize]{L{1.4cm}L{3.9cm}L{4.7cm}L{1.2cm}}
구분 & 자료 & 왜 · 어느 부분 & 시간 \\ \midrule
필수 & Day3 워크북 완성 예시본 5종 & Day4 산출물의 입력 — EX-01 대안 3 · R-14 · 귀책 분해표 · 착시 5 & 60분 \\
필수 & PRINCE2 7 Closing a Project & 5활동 · End Project Report · Follow-on action & 40분 \\
필수 & SRE ch.3 · ch.4 (온라인 공개) & 99.5\% \ra{} 분 \ra{} 릴리스 환전의 원전 & 40분 \\
필수 & 개인정보 보호법 제21445호 조문 & '유출등' 정의 · 72시간 — 직접 읽고 온다 & 30분 \\
심화 & Zwikael \& Smyrk · Ward \& Daniel · Duffield \& Whitty · Williams · Ford 외 · Humble \& Farley ch.10 · Accelerate · ITIL Foundation v5 & 6.1\til{}6.7절의 원전 — 회사 규정 개정 전에 & 25시간 \\
\end{tbl}
\keymsg{읽지 않아도 되는 것 — "종료 체크리스트 100가지"(항목 수가 아니라 세 칸이 문제다) · ITIL 판본 비교 블로그 · 편익 관리 SW 백서}
\note{교재 6.9. 필수 4종은 과정 전 자료이므로 이미 읽었어야 한다 — 읽지 않은 수강생에게는 사후 30일 과제와 함께 다시 안내한다. 심화 8종은 합계 약 25시간이며 순서를 정해 준다: 회사 규정 개정을 맡은 사람은 ITIL Foundation v5부터, 편익 책임을 맡은 사람은 Zwikael \& Smyrk부터, PMO는 Duffield \& Whitty부터. 참고 자료(ISO/IEC 20000-1 조항 · MSP 5판 Benefits · 회사 정본 §5.5·§6.5 · 검산 스크립트 day4\_closing.py)는 워크북을 쓰다 손을 뻗어 확인하는 것이라 화면에 올리지 않았다. 예상 소요 2분.}
\end{frame}

\begin{frame}{4일 20종 산출물의 사슬을 거꾸로 — ⑳에서 ①로}
\fig[0.60]{../그림/PM_Day4/fig_6_1_twenty_chain.pdf}
\figcap{그림 6-1. 편익 원장 6칸 \ra{} 편익 책무 분리(ITO) \ra{} 헌장 §3 성공 기준 \ra{} BC §4 측정 책임자·분모·§6 손익분기 66\%. 여섯 가닥(측정 책임자·분모·손익분기·AC·지체상금 0·PoNR) — 고리가 비면 그 자리에서 숫자가 끊긴다}
\keymsg{20종은 스무 개의 문서가 아니라 하나의 사슬이다 — Day1에 측정 책임자를 부서명으로 적었다면 오늘 ITO 서명표에는 서명할 사람이 없었다}
\note{교재 "PM 트랙을 마치며"·워크북 §7. 그림의 여섯 가닥을 하나씩 손으로 짚는다. ⑳ 원장 "측정 책임자 실명" ← ① BC §4 열 / "분모" ← BC §4 산출 분모 + 가정 A3 / "손익분기 66\%" ← BC §6 / ITO outcome 소유자 ← ② 헌장 §3 + ④ 거버넌스(BC §9). ⑱ 최종 AC 151.95 ← ⑨ PV 곡선(145.80 \ra{} BL-02 146.82) + ⑫ EAC4 150.86 / 지체상금 0 ← ⑧ 여유 소유권 문안 + ⑮ 귀책 분해표 / 이관 6번 ← ⑦ StRS-034 두 분모 리포트 + ① A3. ⑯ PoNR·트리거 ← ⑧ 컷오버 창 허용범위 0 + ⑩ R-14 + ② 헌장 §7. 반사실 세 개를 말한다 — Day2에 두 분모 리포트를 Must로 두지 않았다면 분모 정치는 G5에서 처음 터졌고, Day3에 라이선스 4.70 착시를 걷어 내지 않았다면 최종 AC는 "예상 못 한 10.7억 초과"로 기록되었을 것이다. 예상 소요 5분.}
\end{frame}

\begin{frame}{세 표준의 전환기를 지나는 실무자에게 — 2026년 9월}
\begin{tbl}[\scriptsize]{L{3.0cm}L{4.4cm}L{4.0cm}}
표준 & 무엇이 바뀌었나 & 이 교재의 규율 \\ \midrule
PMBOK 8판 (PMI, 2025) & Focus Area · 성과영역 — 6판 번호와 병기 & 판본을 정직하게 말한다(각 X.2 대조표) \\
ITIL Version 5 (2026-02) & ITIL Value System · Lifecycle 8활동 — practice 명칭 유지 & 어휘를 왜 고정하는지 고지한다("ITIL 4가 최신" 금지) \\
개인정보 보호법 제21445호 (시행 2026-09-11) & '유출등'에 위조·변조·훼손 추가 · 72시간 통지·신고 & 조문은 단정, 적용은 비단정, 시한은 단정 \\
\end{tbl}
\subhead{회사 규정을 개정할 때의 순서}
\begin{itemize}
\item ① 지금 계약서에 적힌 어휘 확인 \ra{} ② 바뀐 상위 구조 매핑 \ra{} ③ 법령의 시한을 절차에 박는다
\end{itemize}
\keymsg{어휘가 낡았다고 규정을 통째로 바꾸는 것과, 판본이 바뀐 줄 모르는 것은 같은 실수의 두 얼굴이다}
\note{교재 "세 표준의 전환기" · 6.7 · 6.8. 이번 주에 개정법이 시행되었다는 사실을 다시 짚는다 — 3.8절의 72시간 타임라인(그림 3-6)은 그 조문 위에 있고, 인쇄 전 조문 원문 대조와 법무 검토는 이 교재가 대신하지 못한다. ITIL v5 사실은 웹 확인한 것이며 8활동 매핑표는 리딩리스트 §1.12에 [검증 필요]로 있다. 규정 개정 순서 셋은 그대로 회사에 가져가도 되는 절차다 — 어휘부터 확인하는 이유는 계약서(SLA·UC)에 적힌 어휘를 바꾸면 계약 변경이 되기 때문이다. 예상 소요 3분.}
\end{frame}

\begin{frame}{Day1\til{}4 — 네 문장}
\begin{tbl}[\small]{L{1.4cm}L{4.2cm}L{5.4cm}}
Day & 한 문장 & 오늘 어디에 남았나 \\ \midrule
Day1 & 숫자 없는 헌장은 소설이다 & 헌장 §3 성공 기준 \ra{} ⑳ ITO 서명표 · BC §6 66\% \\
Day2 & 기준선은 협상된다 & ⑧ 여유 소유권 · ⑨ 145.80 \ra{} 146.82 · ⑩ R-14 \\
Day3 & 편차는 결정 요청으로 끝난다 & ⑫ EX-01 스폰서 결정 \ra{} G4 = 기준선 완료일 \ra{} 지체상금 0 \\
Day4 & 오픈은 이관이다 & ⑯ D+1 07:20 서명 · ⑰ SLA 3층 · ⑱ 이관 10건 · ⑳ 두 시계 \\
\end{tbl}
\keymsg{네 문장은 하나다 — 결정되지 않은 것을 결정하고, 분해해 숫자로 만들고, 편차를 결정으로 되돌리고, 받는 쪽이 서명하게 한다}
\note{Day1\til{}Day3 마무리 교시의 "오늘의 한 문장"을 세로로 세운다. 각 문장이 오늘 어느 산출물에 남았는지를 말하는 것이 이 프레임의 목적이다 — 문장 자체는 이미 들었다. Day3의 문장은 "통제는 편차를 재는 일이 아니라 누가 무엇을 결정해야 하는지를 적는 일"이었고, 오늘 그 결정(EX-01 대안 1, 10-16)이 G4 = 기준선 완료일이 되어 지체상금 0의 첫 줄이 되었다(그림 4-3). 예상 소요 3분.}
\end{frame}

\begin{frame}{Product 트랙 예고 — "당신은 발주자였다. 다음 4일은 공급자다"}
\begin{columns}[T]
\begin{column}{0.48\textwidth}
\subhead{4일 동안 — 발주사 온담의 자리}
\begin{itemize}
\item "대기와 추가는 변경대가 대상"에 WBS 담당 열·RACI·동시 기록으로 답했다
\item 검수 보류 회신 · 감리 귀속 이의 · 지체상금 0과 대기 비용 0.60을 같은 표로
\item 서명한 SAC·SLA·원장 = 인도받은 것을 조직이 받아 쓰기 위한 문서
\end{itemize}
\end{column}
\begin{column}{0.50\textwidth}
\subhead{Day5\til{}8 — 반대편에 선다}
\begin{itemize}
\item 온담은 같은 회사로 다시 등장 — 이번에는 당신의 \textbf{고객}
\item 당신은 한 스타트업의 프로덕트 책임자 — 온담을 \textbf{밖에서} 본다
\item 종료 보고서를 읽을 수 없고, 이관 목록에 무엇이 있는지 모른다
\item 세 본부장의 이음새 · CFO 결재 절차 · 부채비율 187\%를 공급자 관점으로
\end{itemize}
\end{column}
\end{columns}
\keymsg{발주자로 배운 것 — 편익의 분모 · 측정 책임자 실명 · 범위 밖 전제 · 승인 전 항목 · 반려된 견적 — 이 공급자로서 무엇을 만들어 팔 것인가의 재료가 된다}
\note{교재 "Product 트랙 예고" · 워크북 §7. 혼합 기수이므로 Product 트랙만 듣는 수강생이 아니라 PM 트랙 수강생에게 말하는 프레임이다 — 4일 동안의 발주자 경험이 Day5부터 어떻게 뒤집히는지를 예고한다. 핵심은 정보의 비대칭이다: 발주자였을 때 당신은 이관 목록 6번이 무엇인지 알았지만, 공급자로서는 그것을 밖에서 발견해야 한다. 발주자의 사슬을 끝까지 따라가 본 사람만이 공급자로서 그 사슬의 어느 고리가 비어 있는지를 볼 수 있다. Product 트랙의 교재·도구 명칭은 여기서 말하지 않는다 — Day5 도입에서 한다. 예상 소요 3분.}
\end{frame}

\begin{frame}{이관 목록의 한 줄이 출발점 — "승인 전" 네 줄 중 어느 줄이 누군가의 사업이 되는가}
\begin{tbl}[\scriptsize]{L{0.5cm}L{5.6cm}L{1.6cm}L{1.9cm}L{1.7cm}}
\# & 내용 & 소유자 & 기한 & 무엇이 승인돼야 \\ \midrule
3 & 가맹 발주 마감 22:00 운용 전환 — 물류 야간조 편성과 함께 CAB 재상정 & 한동석 & 2027-09 CAB & 야간조 편성안 · 3PL 시각 합의 \\
6 & 이천센터·안성 검사기록은 P1 범위 밖(종이·엑셀). 98.5\% 목표는 디지털화가 전제 — 2025-11 이사회 반려(34억) 재검토 & 나영선 & 2027-09 이사회 & 이사회 재상정 \\
7 & SAP ECC 2027-12 유지보수 종료 — 애드온·IF-SAP 8본 마이그레이션 계획(KE-04) & 박세연 & 2027-08-31 & 별도 과제 예산 \\
10 & 하자보수기 PMO 잔류 1명 — 유보 3.0 중 1.05 집행 & 프로그램 PMO장 & 2027-06-01 운영위 & 스폰서 결정 \\
\end{tbl}
\keymsg{프로젝트가 끝난 뒤에도 남아 작동하는 문서는 그 프로젝트의 것만이 아니다 — 다음 프로젝트의, 다음 사람의 출발점이다. Day5에 만나자}
\note{워크북 §3.5 항목 6(이관 10건 중 승인 전 4)과 §7의 마지막 질문. 네 줄을 읽고 수강생에게 1분 생각할 시간을 준다 — 어느 줄이 "회사 안의 후속 과제"로 끝나고, 어느 줄이 "밖의 누군가가 만들어 팔 물건"이 되는가. 답을 정해 주지 않는다. 6번(그림 4-4 강조 · 그림 5-3의 분모 B · BR-05)이 편익 98.5\%의 전제이면서 34억 견적이 반려된 자리라는 점만 짚는다 — 반려된 견적은 문제가 사라졌다는 뜻이 아니라 그 값으로는 안 산다는 뜻이다. 예상 소요 3분.}
\end{frame}

\begin{frame}{사후 읽기 · 30일 과제 — 당신 회사의 종료 보고서 한 부로}
\begin{itemize}
\item \textbf{사후 읽기(2주)} — PRINCE2 7 Closing a Project · SRE ch.3\til{}4 · 개정법 조문(다시) · Duffield \& Whitty
\item \textbf{과제 ① 세 칸 O/X} — 원가는 어느 층과 비교됐나 / 이관 목록에 소유자 실명·승인 여부가 있나 / 측정 책임자가 실명인가
\item \textbf{과제 ② 거꾸로 추적} — 그 보고서의 편익 한 줄을 BC·헌장까지 거슬러 고리가 비는 지점 하나를 적는다
\item \textbf{과제 ③ 한 장} — 비는 고리 하나를 채우는 개정 요청(교훈 등록부 ⑤ 소유자 · Syllk 층)을 A4 한 장으로
\item 제출: Day8 종료 전 · 형식은 워크북 §4·§5의 양식 그대로
\end{itemize}
\keymsg{항목의 수가 아니라 세 칸이 문제다 — 30일 뒤 당신 회사의 종료 보고서에 그 세 칸이 생기면 이 4일은 끝난 것이다}
\note{교재 6.9의 과정 전 과제(세 칸 O/X)를 사후 과제로 되돌린다 — 과정 전에는 "읽고 O/X만", 사후에는 "비는 고리를 채우는 개정 요청 한 장"까지. 과제 ②는 그림 6-1의 방법을 자기 회사 문서에 적용하는 것이며, 고리가 비는 지점은 대개 측정 책임자(부서명)·분모(미정의)·승인 전 항목(승인처럼 적힘) 셋 중 하나다. 과제 ③은 워크북 §4 교훈 등록부의 ④ 적용 대상·⑤ 소유자 칸과 §4의 개정 요청 양식을 그대로 쓴다 — 양식이 있어야 process·infrastructure 층(Syllk)에 들어간다. 제출 시점을 Day8 종료 전으로 두어 Product 트랙 수강 중에도 PM 트랙의 문서가 살아 있게 한다. 예상 소요 3분. 이로써 마무리 30분 — 2+4+4+2+5+3+3+3+3+3 = 32분, 필요시 읽을 자료 프레임을 1분으로 줄인다.}
\end{frame}
